% ---
% id: meat-core/cor:separate-executions
% kind: cor
% title: "Security parameters of separate executions add"
% summary: >-
%   N executions that share no key and no device state and leave their keys
%   unused until all have ended, run against one adversary, are together
%   (ε_1 + ... + ε_N)-secure for the joint ideal map. A key reused from one
%   execution in another is not covered.
% user_statement: >-
%   Running several executions against the same attacker, one after another
%   or at the same time, costs the sum of their security parameters. If each
%   execution i is secure with parameter ε_i and [[A-SEPARATE]] holds, whose
%   conditions include that their number and parameters are fixed before the
%   first starts, that they share no key material and no device state, and
%   that their keys stay unused until all have ended, then all their keys
%   together are within ε_1 + ... + ε_N of independent ideal keys. Parameters
%   chosen from an earlier execution's data are not covered, nor is a key
%   from one execution reused in another, for example to authenticate it.
% prominence: standard
% status: accepted
% uses: [cor:composed-security, thm:keylength, thm:accumulated-entropy, thm:constrained-accumulation, lem:ev-correctness, def:roundmodel, def:execution-end]
% mentions: [protocol-compiler/ass:atk]
% assumes: [A-SEPARATE]
% refs:
%   - "tupkary2026decoy, Sec. 3.3, Def. 3.8, Eqs. (13)-(17)"
%   - "portmann2022security, Sec. II.E"
% provenance: newly-derived
% lean:
% ---

Composable frameworks add the errors of protocols composed in sequence or
in parallel \cite[Sec.~II.E]{portmann2022security}; that is the motivation
for the trace-distance statement below, which is derived without it. Each
execution's parameter holds against every attack on it, and the other
executions, which touch none of its registers, are part of one such
attack. The per-execution ideal maps are then applied one at a time, each
step costing that execution's parameter.

\begin{corollary}\label{cor:separate-executions}
Let executions $1,\dots,N$ of the protocol run against one adversary, in
sequence, concurrently or interleaved, the adversary carrying information
from each to the others. For each $i$ let the hypotheses of
Corollary~\ref{cor:composed-security}\,(iii) hold for execution~$i$, with
its parameters $\epa^{(i)}$, $\eev^{(i)}$ and
$\varepsilon_{\mathrm{auth}}^{(i)}$. Write
$K^{(i)}\coloneqq K_A^{(i)}K_B^{(i)}$ for its key registers,
$\mathcal R^{(i)}$ for the map $\mathcal R_{\mathrm{ideal}}$ of
\cite[Sec.~3.3, Eqs.~(13)--(17)]{tupkary2026decoy} on $K^{(i)}$, and
$\varepsilon_i\coloneqq\epa^{(i)}+\eev^{(i)}+2\varepsilon_{\mathrm{auth}}^{(i)}$.
Let Assumption~\ref{ass:separate-executions} hold. Then for every attack on
the $N$ executions under which each of them reaches its
\tref{end of an execution} (Definition~\ref{def:execution-end}) on every
branch, with $\rho$ the joint state of $K^{(1)}\cdots K^{(N)}$ and the
adversary's information $E$ at the end of the last of them,
\begin{equation}\label{eq:separate-executions}
  \frac12\Big\|\rho-\big(\mathcal R^{(1)}\otimes\cdots\otimes
  \mathcal R^{(N)}\otimes\idmap_E\big)(\rho)\Big\|_1
  \le\sum_{i=1}^{N}\varepsilon_i .
\end{equation}
\end{corollary}

\begin{proof}
Fix an attack on the $N$ executions under which each of them reaches its
end on every branch. For each $i$ let $t_i$ be the
\tref{end of an execution} (Definition~\ref{def:execution-end}) for
execution~$i$, and let
$t\coloneqq\max_i t_i$. For a time $s$ and an execution $i$, write
$Y_i(s)$ for the adversary's registers at $s$ together with every register
of the honest parties of every execution $j\ne i$ at $s$. By
Assumption~\ref{ass:separate-executions}\,(iii), at $t$ the register
$K^{(j)}$ of every execution holds that execution's output.

The times $t_i$ and $t$ are random and may depend on the attack. The joint
process is therefore read branch by branch over the classical record of
when each party ends (where the end times are bounded, equivalently with
idle steps that give every branch one schedule); the states named below
are averaged over the branches. That record lies in $Y_i$: the key holder of execution~$i$ ends
by sending a message the adversary sees, and the partner when the
adversary delivers it (Definition~\ref{def:execution-end}). By
Assumption~\ref{ass:separate-executions}\,(ii) and (iii), the evolution
after execution~$i$ has ended depends on that execution's history only
through $Y_i$. The restriction to attacks under which each execution
reaches its end on every branch is the conjunction of the restrictions
that Corollary~\ref{cor:composed-security}\,(iii) already makes for each
execution, whose bound concerns the output state of a run that ends.

\emph{Step 1: one execution against the others.} Fix $i$. Let
$\mathcal E_i$ be the adversary whose initial register holds the given
adversary's initial register together with the initial registers of the
honest parties of every execution $j\ne i$ (their pre-shared keys, their
local randomness and the initial states of their devices), in their joint
initial state, and who carries out, at the times and in the order of the
given attack, the given adversary's operations and every operation of the
honest parties of the executions $j\ne i$. By
Assumption~\ref{ass:separate-executions}\,(ii), no operation of the honest
parties of another execution acts on a register of execution~$i$'s honest
parties, none of theirs acts on a register of another execution, and
execution~$i$ meets the other executions only through systems and
classical channels the adversary has access to. The given adversary's own
operations reach execution~$i$ only through those systems and channels,
its attack interface, as in any attack on it. So
$\mathcal E_i$ is an attack on execution~$i$ alone, whose registers at time
$s$ are $Y_i(s)$, and execution~$i$ run against $\mathcal E_i$ is the same
channel, applied to the same initial state, as the $N$ executions against
the given adversary, with the registers grouped differently.

Corollary~\ref{cor:composed-security}\,(iii) holds for every attack that
its premises admit, whatever the dimension of the adversary's registers, and
it restricts the adversary's initial register and its state by nothing
else. Through Theorem~\ref{thm:keylength} it rests on
Theorem~\ref{thm:accumulated-entropy}, whose initial state $\omega^0$ is
arbitrary on the adversary's initial register $E_0$ of the
\tref{round model} (Definition~\ref{def:roundmodel}), with Alice's marginal
fixed in \tref{trusted-source mode}; in the constrained branch of
Theorem~\ref{thm:keylength} it rests on
Theorem~\ref{thm:constrained-accumulation}, whose hypotheses include those
of Theorem~\ref{thm:accumulated-entropy} and so leave $\omega^0$ as
arbitrary, the \tref{attack sequence} being $\CList$-admissible. It also
rests on Lemma~\ref{lem:ev-correctness}, which holds for every attack. The
initial register of $\mathcal E_i$, however it
is correlated with the given adversary's and the other executions'
systems, is such an $E_0$. The premises are the assumptions on which
Corollary~\ref{cor:composed-security}\,(iii) rests for execution~$i$, and
by Assumption~\ref{ass:separate-executions}\,(iv) each of them holds for
execution~$i$ against $\mathcal E_i$. The registers of $\mathcal E_i$,
which hold the other executions' registers, may be of any dimension, which
the corollary allows. Where a constraint list is declared, A-ATK
(\xref{protocol-compiler/ass:atk}) states its reduction over the
\tref[licensed attack set]{licensed attack sets} of
Definition~\ref{def:roundmodel}, defined for finite-dimensional registers,
and is asked, with the corollary's other premises, of the all-at-once attacks
with finite registers through which the corollary reaches $\mathcal E_i$.
Clause (iv) asserts each of these, and the proof uses them only through
Corollary~\ref{cor:composed-security}\,(iii). By
Assumption~\ref{ass:separate-executions}\,(i) the parameters of
execution~$i$, and with them $\varepsilon_i$, are fixed before any
execution starts and do not depend on the attack.

Corollary~\ref{cor:composed-security}\,(iii) therefore applies to
execution~$i$ against $\mathcal E_i$, whose information when execution~$i$
has ended is $Y_i(t_i)$. With $\sigma_i$ the state of $K^{(i)}Y_i(t_i)$ at
$t_i$, \eqref{eq:composed-auth} gives
\begin{equation*}
  \frac12\big\|\sigma_i-(\mathcal R^{(i)}\otimes\idmap)(\sigma_i)\big\|_1
  \le\varepsilon_i .
\end{equation*}
From $t_i$ to $t$, nothing acts on $K^{(i)}$ and no information about
another register of execution~$i$'s honest parties leaves their
laboratories (Assumption~\ref{ass:separate-executions}\,(iii)); those
parties, having ended, act on nothing; and no operation of the honest
parties of another execution acts on a register of execution~$i$
(Assumption~\ref{ass:separate-executions}\,(ii)). What happens between
$t_i$ and $t$ therefore acts on $Y_i$ alone and takes no input from a
register of execution~$i$: it is a channel $\Phi_i$ from $Y_i(t_i)$ to
$Y_i(t)$, and the state of $K^{(i)}Y_i(t)$ at $t$ is
$(\idmap_{K^{(i)}}\otimes\Phi_i)(\sigma_i)$. The registers $K^{(j)}$,
$j\ne i$, and $E$ at $t$ belong to $Y_i(t)$, so
$\rho=(\idmap_{K^{(i)}}\otimes\Lambda_i)(\sigma_i)$ with $\Lambda_i$ the
channel $\Phi_i$ followed by the partial trace onto them. This channel acts
on registers other than $K^{(i)}$, so it commutes with
$\mathcal R^{(i)}\otimes\idmap$, and a channel does not increase the trace
norm of a Hermitian operator. Hence
\begin{equation}\label{eq:separate-one}
  \frac12\big\|\rho-(\mathcal R^{(i)}\otimes\idmap)(\rho)\big\|_1
  \le\varepsilon_i\qquad(i=1,\dots,N).
\end{equation}

\emph{Step 2: the hybrid.} Put $\mathcal T_0\coloneqq\idmap$ and
$\mathcal T_i\coloneqq\mathcal R^{(1)}\otimes\cdots\otimes
\mathcal R^{(i)}\otimes\idmap$ for $i=1,\dots,N$, the identity acting on
every other register. The maps $\mathcal R^{(i)}$ act on the disjoint
registers $K^{(i)}$, so
$\mathcal T_i=\mathcal T_{i-1}\circ(\mathcal R^{(i)}\otimes\idmap)$, and
$\mathcal T_N$ is the map of \eqref{eq:separate-executions}. By linearity,
\[
  \rho-\mathcal T_N(\rho)
  =\sum_{i=1}^{N}\big(\mathcal T_{i-1}(\rho)-\mathcal T_i(\rho)\big)
  =\sum_{i=1}^{N}\mathcal T_{i-1}\big(\rho-(\mathcal R^{(i)}\otimes
  \idmap)(\rho)\big).
\]
Each $\mathcal R^{(i)}$ is a channel
\cite[Sec.~3.3, Eqs.~(13), (17)]{tupkary2026decoy}, so each
$\mathcal T_{i-1}$ is a channel and does not increase the trace norm of the
Hermitian operator it is applied to. The triangle inequality and
\eqref{eq:separate-one} give \eqref{eq:separate-executions}.

By \cite[Sec.~3.3, Eqs.~(14), (16), (17)]{tupkary2026decoy},
$\mathcal R^{(i)}$ replaces, for each pair of key lengths of execution~$i$,
its key registers by the ideal keys of that pair, independent of every other
register. Applied for one execution after another, $\mathcal T_N(\rho)$ is
therefore the state in which, for each tuple of the executions' pairs of key
lengths, with its probability under $\rho$, the key registers of every
execution hold the ideal keys of its pair, independent of each other and of
$E$: the ideal of the $N$ executions together.
\end{proof}
